\documentclass[12pt]{article}%
\usepackage{amsfonts}
\usepackage{fancyhdr}
\usepackage{comment}
\usepackage[a4paper, top=2.5cm, bottom=2.5cm, left=2.2cm, right=2.2cm]%
{geometry}
\usepackage{times}
\usepackage{amsmath}
\usepackage{changepage}
\usepackage{amssymb}
\usepackage{graphicx}

\setcounter{MaxMatrixCols}{30}
\newtheorem{theorem}{Theorem}
\newtheorem{acknowledgement}[theorem]{Acknowledgement}
\newtheorem{algorithm}[theorem]{Algorithm}
\newtheorem{axiom}{Axiom}
\newtheorem{case}[theorem]{Case}
\newtheorem{claim}[theorem]{Claim}
\newtheorem{conclusion}[theorem]{Conclusion}
\newtheorem{condition}[theorem]{Condition}
\newtheorem{conjecture}[theorem]{Conjecture}
\newtheorem{corollary}[theorem]{Corollary}
\newtheorem{criterion}[theorem]{Criterion}
\newtheorem{definition}[theorem]{Definition}
\newtheorem{example}[theorem]{Example}
\newtheorem{exercise}[theorem]{Exercise}
\newtheorem{lemma}[theorem]{Lemma}
\newtheorem{notation}[theorem]{Notation}
\newtheorem{problem}[theorem]{Problem}
\newtheorem{proposition}[theorem]{Proposition}
\newtheorem{remark}[theorem]{Remark}
\newtheorem{solution}[theorem]{Solution}
\newtheorem{summary}[theorem]{Summary}
\newenvironment{proof}[1][Proof]{\textbf{#1.} }{\ \rule{0.5em}{0.5em}}

\newcommand{\Q}{\mathbb{Q}}
\newcommand{\R}{\mathbb{R}}
\newcommand{\C}{\mathbb{C}}
\newcommand{\Z}{\mathbb{Z}}

%%%%
% ME %
%%%%

\usepackage{listings}
\usepackage{color}
\usepackage{xcolor}
\usepackage{mdframed}

\definecolor{dkgreen}{rgb}{0,0.6,0}
\definecolor{gray}{rgb}{0.5,0.5,0.5}
\definecolor{mauve}{rgb}{0.58,0,0.82}

\lstset{%frame=tb,
language=Matlab,
aboveskip=3mm,
belowskip=3mm,
showstringspaces=false,
columns=flexible,
basicstyle={\small\ttfamily},
numbers=none,
numberstyle=\tiny\color{gray},
keywordstyle=\color{blue},
commentstyle=\color{dkgreen},
stringstyle=\color{mauve},
breaklines=true,
breakatwhitespace=true,
tabsize=3
}

\usepackage{outlines}
\usepackage{enumitem}
%\setenumerate[1]{label=\Roman*.}
%\setenumerate[2]{label=\Alph*.}
%\setenumerate[3]{label=\roman*.}
%\setenumerate[4]{label=\alph*.}

\usepackage{titlesec} 

\titleformat{\subsection}[runin]
  {\normalfont\large\bfseries}{\thesubsection}{1em}{}
\titleformat{\subsubsection}[runin]
  {\normalfont\normalsize\bfseries}{\thesubsubsection}{1em}{}

% for enumerate spacing
\newenvironment{tight_enum}{
\begin{enumerate}[label=\alph*.]
\setlength{\itemsep}{0pt}
\setlength{\parskip}{0pt}
}{\end{enumerate}}

\usepackage{exercise}

\setlength\parindent{0pt}

\newcommand{\code}[1]{\texttt{#1}}
\newcommand{\shp}{$>>>$ }
\newcommand{\tab}{\hspace*{45pt}}

\begin{document}

\title{Homework 03: Chapter 04 and 06}

\author{EEC 3150}
\date{}

\maketitle

\begin{center}
\fbox{\fbox{\parbox{5.5in}{\centering
For each Python exercise, write a separate Python script. Submit all scripts on Blackboard under the appropriate assignment. Name your script files with the following format: \\
\vspace{10pt}
HW$<$assignment number$>$\_Ex$<$exercise number$>$\_$<$FirstInitial$><$LastName$>$.py \\
\vspace{10pt}
Example: HW01\_Ex01\_GWest.py \\
\vspace{10pt}
For short answer or math exercises, either do your work on a sheet of paper and upload a picture of it to Blackboard or submit some kind of text file or Word document to Blackboard. You can place all such exercises in a single document as long as you clearly label each exercise. These files should follow a similar naming convention as the Python scripts.
}}}
\end{center}


\pagebreak


\begin{Exercise}[origin={10 points, Python}]

Write a function \code{Reverse(s)} which will accept a single string argument and return the reversed string. Call the function. Here's an example of how to use the function:

\vspace{10pt}
\code{\shp Reverse(`abcdefg') \\
'gfedcba'
}

\vspace{10pt}
DO NOT use any built-in string methods.

\end{Exercise}

\begin{Exercise}[origin={10 points, Python}]

Write a function \code{Split(s,n)} which will accept two arguments (a string and an integer $n$) and return two values (the left and right halves of the string when split at the $n$-th index of the string). Call the function. Here's an example of how to use the function:

\vspace{10pt}
\code{\shp Split(`abcd', 2) \\
(`ab', `cd')
}

\vspace{10pt}
DO NOT use any built-in string methods.

\end{Exercise}

\begin{Exercise}[origin={10 points, Python}]

Write a function \code{ConcatStrings(*s)} which will accept a variable number of string arguments and return the concatenation of all of them. Demonstrate the function by calling it. Here's an example of how to use the function:

\vspace{10pt}
\code{\shp ConcatStrings(`a','b','c') \\
'abc'
}

\end{Exercise}

\begin{Exercise}[origin={10 points, Python}]

Write a function \code{trig(x,n)} which will accept two arguments (a float $x$ and an integer $n$) and return two values: the $n$-th order Taylor approximation of sin($x$) and cos($x$).
\begin{equation*}
\begin{split}
sin(x) & \approx  \sum_{i = 0}^n \frac{(-1)^n x^{2n+1}}{(2n+1)!} \\
cos(x) & \approx  \sum_{i = 0}^n \frac{(-1)^n x^{2n}}{(2n)!}
\end{split}
\end{equation*}
Test the function by plugging in an approximation for $\pi$ for $x$. You may use my factorial code or write your own.

\end{Exercise}

\begin{Exercise}[origin={20 points, Python}]

Write a program with two functions---one non-recursive (\code{HarmNon(n)}) and one recursive (\code{HarmRec(n)})---that will calculate the Harmonic series:
\begin{equation*}
	\sum_{i=1}^{n} \frac{1}{i} = \frac{1}{1} + \frac{1}{2} + \cdots + \frac{1}{n}
\end{equation*}

Test your function by printing the values of both functions with something like the following code:

\vspace{10pt}
\code{n = 10 \\
for i in range(1,n+1): \\
\tab print(HarmNon(i),HarmRec(i))
}
\end{Exercise}




\end{document}


